\section{Experimental Setup}
\label{sec:setup}

\paragraph{Methodology and provenance.}
The evaluation has five layers, each reproducible from \texttt{configs/main.yaml}:
(1) the \emph{closed-form model} of \S\ref{sec:design}--\S\ref{sec:reconf};
(2) \emph{Monte-Carlo simulators} that independently estimate the two probabilistic quantities
($\psucc$ from the BFT quorum harness, $\pwin$ from the accountability-race simulator) and are checked
against the closed form;
(3) a \emph{local proof-of-concept} that runs the end-to-end attack on an abstract BFT committee,
producing real equivocation evidence and a paying escrow;
(4) a \emph{discrete-event testbed} with real Ed25519 signatures in which $\Tacc$ and $\Tsettle$ are
measured from a running protocol and cross-validated against the closed form; and
(5) a \emph{case study} that evaluates the model at the documented parameters of deployed systems. No
quantity in \S\ref{sec:eval} is a hand-chosen placeholder: each is computed by the artifact, and the
figures/tables are regenerated by the commands of \S\ref{sec:impl}. We do not \emph{actively} measure
a live deployed network under attack; that is the stated next step in \S\ref{sec:limitations}, and the
testbed and case study are structured so that each quantity has a live counterpart to instrument.

\paragraph{Live-measurement hook.}
For a concrete deployment, the two load-bearing quantities can be measured without executing an
attack by instrumenting timestamps for: source certificate creation; destination receipt acceptance;
withdrawal or bridge finalization; first observation of conflicting certificates; evidence
propagation to the accountability layer; slashing/quarantine effectiveness; and stake withdrawal or
unbonding completion. These logs give $\Tsettle$ directly, give $\Tacc$ from evidence observation to
effective enforcement, and estimate
$q=\Pr[t_{\mathrm{slash}}<t_{\mathrm{withdraw}}]$ under the deployment's normal evidence and
unbonding rules. The local testbed uses the same timestamp schema, which is why its outputs are a
drop-in rehearsal for the production measurement.

\paragraph{Parameters and sweeps.}
Table~\ref{tab:params} lists parameters and swept ranges. We study committee sizes from small
practical committees to large ones; enforcement probability $q$ from ``accountability rarely lands in
time'' to ``near-certain''; and the latency ratio $\rho=\Tacc/\Tsettle$ across the critical point
$\rho=1$. We treat $\psucc$ as the adversary's split-control reliability; we fix it at the strong-split
value $0.85$---the \emph{view-control} regime, or $k\approx F{+}6$ at ${\approx}1.45\times$ the floor
bribe by an uncontrolled split (\S\ref{sec:floor})---for the window analyses, and sweep it
$\psucc\in[\ArtifactPsuccFloor,1]$ for sensitivity, reporting alongside the random-split \emph{floor}
of Eq.~\eqref{eq:psucc}.

\begin{table}[t]
\centering
\caption{Simulation parameters and swept ranges (from \texttt{configs/main.yaml}).}
\label{tab:params}
\small
\begin{tabularx}{\linewidth}{@{}llX@{}}
\toprule
Parameter & Symbol & Range (swept) \\
\midrule
Committee size & $N=3F{+}1$ & $\{4,7,\ldots,400\}$ (8 sizes) \\
Equivocators & $k=F{+}1$ & derived from $N$ \\
View-split & $\varphi$ & $[0.02,\,0.98]$ \\
Stake at risk & $s$ & normalized to $1$ \\
Continuation value & $R$ & $\{0,\,0.5\,s,\,s\}$ \\
Enforcement prob. & $q$ & $[0.05,\,1.0]$ \\
Latency ratio & $\rho=\Tacc/\Tsettle$ & $[0.25,\,4]$ \\
Settle-time CV & $\mathrm{CV}_{\Tsettle}$ & $\{0.1,0.3\}$ \\
Extractable value & $\Eval$ & $[0,\,200]\cdot s$ \\
Split-control & $\psucc$ & floor $\ArtifactPsuccFloor$ to $1$; $0.85$ = view-control \\
Epoch length & $\tau$ & $1$ (normalized) \\
Unbonding delay & $U$ & $0.15\,\tau$ \\
Cross-epoch latencies & $\Treceipt,\Tstate$ & $0.10,\,0.08$ \\
Handoff penalty (max) & $\Hreconf$ & $0.40$ \\
Monte-Carlo trials & --- & up to $6{\times}10^4$ \\
\bottomrule
\end{tabularx}
\end{table}

\paragraph{Metrics.}
We report (i) $\psucc$ vs.\ $\varphi$ for $k\in\{F{+}1,F{+}2,F{+}3\}$ with a Monte-Carlo overlay;
(ii) the window-open probability $\pwin$ vs.\ $\rho$ and settle-time variance; (iii) the minimum bribe
$B$ vs.\ $N$ and $q$; (iv) the break-even value $\Eval^{*}=B+\Cop$; (v) expected profit
$\mathbb{E}[\Pi]$ over the $(\rho,\Eval)$ and $(q,\Eval)$ planes; (vi) $\mathbb{E}[\Pi(a)]$ across the
epoch for each handoff design; and (vii) the effect of each defense on $\mathbb{E}[\Pi]$.

\paragraph{Baselines.}
As an immunity baseline we include the finality-gated configuration of
Proposition~\ref{prop:immunity} ($\pwin\equiv0$), against which any profitable region must vanish, and
the accountability-gated carryover of Proposition~\ref{prop:reconf-safe} for the reconfiguration
sweep.
